% Part: computability
% Chapter: computability-theory
% Section: halting-problem

\documentclass[../../../include/open-logic-section]{subfiles}

\begin{document}

\olfileid{cmp}{thy}{hlt}
\olsection{مسئلهٔ توقف}

بنا بر ساخت، تابع محاسبه‌پذیر جزئیِ جهانی~$\fn{Un}(e,x)$ تعریف شده
است اگر و تنها اگر محاسبهٔ تابعی که~$e$ آن را کد می‌کند برای ورودی
$x$ مقداری تولید کند. طبیعی است بپرسیم آیا می‌توان دربارهٔ وقوع این
حالت تصمیم گرفت. در واقع نمی‌توان. در مدل محاسبهٔ ماشین تورینگ، این
یعنی از نظر محاسباتی تصمیم‌ناپذیر است که آیا ماشین تورینگ داده‌شده‌ای
روی ورودی داده‌شده متوقف می‌شود یا نه. ازاین‌رو قضیهٔ زیر به
«تصمیم‌ناپذیری مسئلهٔ توقف» معروف است. در ادامه دو برهان می‌آوریم.
نخستین برهان رشتهٔ بحث پیشین را دنبال می‌کند و دومی مستقیم‌تر است.

\begin{thm}
\ollabel{thm:halting-problem}
بگذارید
\[
h(e, x)  =
\begin{cases}
1 & \text{اگر $\fn{Un}(e, x)$ تعریف شده باشد} \\
0 & \text{در غیر این صورت.}
\end{cases}
\]
آنگاه $h$ محاسبه‌پذیر نیست.
\end{thm}

\begin{proof}
فرض کنید $h$ محاسبه‌پذیر باشد. نشان می‌دهیم این فرض امکان تعریف یک
تابع محاسبه‌پذیر جهانی را می‌دهد. تعریف کنید:
\[
\fn{Un'}(e,x) =
\begin{cases}
\fn{Un}(e,x) & \text{اگر $h(e,x) = 1$} \\
0 & \text{در غیر این صورت.}
\end{cases}
\]
اکنون $\fn{Un'}(e,x)$ تابعی تام است و اگر~$h$ محاسبه‌پذیر باشد، آن
نیز محاسبه‌پذیر است. برای نمونه، می‌توانیم $g$ را با بازگشت اولیه
چنین تعریف کنیم:
\begin{align*}
g(0, e, x) & \simeq 0 \\
g(y+1, e, x) & \simeq \fn{Un}(e,x);
\end{align*}
سپس
\[
\fn{Un'}(e,x) \simeq g(h(e,x),e,x).
\]
چون هر جا $\fn{Un}(e,x)$ تعریف شده باشد، $\fn{Un'}(e,x)$ با آن
برابر است، $\fn{Un'}$ برای آن دسته از توابع محاسبه‌پذیر جزئی که
اتفاقاً تام‌اند جهانی است. اما این با
\olref[nou]{thm:no-univ} تناقض دارد.
\end{proof}

\begin{proof}
فرض کنید $h(e,x)$ محاسبه‌پذیر باشد. تابع $g$ را چنین تعریف کنید:
\[
g(x) =
\begin{cases}
  0 & \text{اگر $h(x,x) = 0$} \\
  \fundefined & \text{در غیر این صورت.}
\end{cases}
\]
تابع $g$ محاسبه‌پذیر جزئی است. برای نمونه، می‌توان آن را به‌صورت
$\umin{y}{h(x,x)=0}$ تعریف کرد. پس برای یک~$e$، به ازای هر~$x$ داریم
$g(x)\simeq\fn{Un}(e,x)$. آیا $g$ در~$e$ تعریف شده است؟ اگر چنین
باشد، بنا بر تعریف~$g$، $h(e,e)=0$ (زیرا $g$ تنها هنگامی تعریف شده
است که $h(e,e)=0$ باشد). بنا بر تعریف~$h$، این یعنی
$\fn{Un}(e,e)$ تعریف‌نشده است. با فرض اینکه به ازای هر~$x$،
$g(x)\simeq\fn{Un}(e,x)$، نتیجه می‌شود $g(e)$ تعریف‌نشده است که
تناقض دارد. از سوی دیگر، اگر $g(e)$ تعریف‌نشده باشد، آنگاه
$h(e,e)\neq0$ و پس $h(e,e)=1$. در نتیجه $\fn{Un}(e,e)$ تعریف شده
است. اما چون $g(x)\simeq\fn{Un}(e,x)$، آنگاه $g(e)$ نیز باید تعریف
شده باشد. باز هم به تناقض می‌رسیم.
\end{proof}

\begin{tagblock}{TMs}
\begin{explain}
می‌توانیم این استدلال را برحسب ماشین‌های تورینگ وصف کنیم. فرض کنید
ماشین تورینگی چون~$H$ وجود داشته باشد که توصیف ماشین تورینگ~$E$ و
ورودی~$x$ را بگیرد و تصمیم بگیرد آیا $E$ روی ورودی~$x$ متوقف می‌شود
یا نه. آنگاه می‌توانستیم ماشین تورینگ دیگری چون~$G$ بسازیم که یک
ورودی~$x$ می‌گیرد، $H$ را اجرا می‌کند تا تصمیم بگیرد آیا ماشین
$M_x$ با نمایهٔ~$x$ روی ورودی~$x$ متوقف می‌شود، و سپس خلاف آن را
انجام می‌دهد. به بیان دیگر، اگر $H$ گزارش دهد که $M_x$ روی ورودی~$x$
متوقف می‌شود، $G$ وارد حلقه‌ای بی‌نهایت می‌شود؛ و اگر $H$ گزارش دهد
که $M_x$ روی ورودی~$x$ متوقف نمی‌شود، $G$ صرفاً متوقف می‌شود. آیا
$G$ هنگامی که نمایهٔ خودش ورودی آن است متوقف می‌شود؟ استدلال بالا
نشان می‌دهد که متوقف می‌شود اگر و تنها اگر متوقف نشود---تناقض. پس
فرض وجود چنین ماشین تورینگ~$H$ای باید نادرست باشد.
\end{explain}
\end{tagblock}

\end{document}
